\section{Related Work}
\label{sec:related}

\subsection{Agricultural Machinery Trajectory Analysis}

GNSS-based trajectory analysis is the dominant foundation for agricultural machinery monitoring. Early and rule-based methods commonly used speed, heading, and field-boundary constraints to distinguish field and road segments. More recent learning-based work improves robustness by modeling temporal structure and trajectory-derived features. Zhai et al.~\cite{zhai2024bilstm} studied sparse agricultural machinery trajectories and used recurrent models to improve trajectory reconstruction and downstream classification. Zhai et al.~\cite{zhai2024atrnet} transformed trajectory sequences into image-like representations for machinery mode classification. Pan et al.~\cite{pan2024chaos} optimized field-road segmentation parameters with a chaos sensing slime mould algorithm.

These methods demonstrate that trajectory time series contain strong operational cues. However, they also expose the limitation that motivates this paper: motion alone cannot fully determine machine state. Classes that differ in engine status, crop intake, unloading, or header engagement may occupy nearly identical regions in the trajectory feature space. Eleven-class \task{} therefore requires additional sensory evidence.

\subsection{Visual and Audio Cues for Machinery Operation}

Visual observations provide scene-level and action-level context that is unavailable in GNSS trajectories. Images can indicate crop density, field boundary, road surface, header position, and whether the machine is actively interacting with crops. Vision transformers pretrained on large-scale image corpora~\cite{dosovitskiy2021vit} are well suited for this modality, especially when adapted to agricultural scenes with lightweight trainable modules.

Audio provides a complementary view of machine process state. Engine-off waiting, idling, and unloading can be difficult to separate spatially, but they differ acoustically. Prior work on ComPASS~\cite{guo2025compass} showed that self-supervised acoustic representation learning is useful for agricultural machinery trajectory time-series classification. The Audio Spectrogram Transformer (AST)~\cite{gong2021ast} further provides a strong pretrained backbone for log-mel spectrogram classification.

\subsection{Multimodal Fusion}

Multimodal learning methods are commonly grouped into early fusion, late fusion, and intermediate fusion~\cite{baltrusaitis2018multimodal}. Early fusion concatenates low-level or encoder-level features; late fusion combines prediction scores from independent models; intermediate fusion enables cross-modal interaction through attention or other learned communication modules. Attention bottleneck models~\cite{nagrani2021attention} show that compact shared tokens can mediate information exchange without full quadratic cross-modal attention.

The present paper uses direct concatenation as a controlled trimodal reference, but adopts class-adaptive logit fusion as the current multimodal model. This design keeps the reproducibility of a late-fusion baseline while allowing the final decision for each operation class to emphasize the modality that carries the most discriminative evidence.
